\documentclass[11pt]{article}

\usepackage[left=0.8in,right=0.8in,top=0.75in,bottom=0.75in]{geometry}
\usepackage[T1]{fontenc}
\usepackage[utf8]{inputenc}
\usepackage[none]{hyphenat}
\usepackage{setspace}

\pagestyle{empty}
\setlength{\parskip}{0.7em}
\setlength{\parindent}{0pt}

\begin{document}

\begin{flushleft}
Abin Daniel Zorto\\
Intelligent Technologies Research Group\\
University of East London\\
University Way, London E16 2RD, UK\\
\texttt{u2091940@uel.ac.uk}
\end{flushleft}

\vspace{0.8em}

\begin{flushleft}
Professor Jie Lu\\
Editor in Chief\\
\emph{Knowledge-Based Systems}\\
Elsevier
\end{flushleft}

\vspace{0.8em}

\begin{flushleft}
\today
\end{flushleft}

\vspace{0.8em}

Dear Professor Lu,

We wish to submit our original research article, \emph{Sparse Biomarker Discovery from Low-Montage Resting-State EEG for tDCS Treatment Response in Adults with Treatment-Resistant Depression}, for consideration for publication in \emph{Knowledge-Based Systems}.

This study addresses patient-level prediction of response to transcranial direct current stimulation from low-montage resting-state EEG in adults with treatment-resistant depression. Using data collected in a remote clinical trial setting, we evaluated a controlled sweep of time windows and sparse feature-selection budgets across a hybrid deep model and a linear baseline under leave-one-participant-out validation, with class rebalancing and feature selection restricted to the training participants within each fold.

The main finding is that strong discriminatory ability was achieved under an extremely sparse feature budget. The best hybrid configuration reached a patient-level ROC-AUC of 0.816, balanced accuracy of 0.786, and MCC of 0.571 across 21 held-out participants while selecting only one feature per fold and seven recurrent features overall. The recurrent biomarkers centered on temporal spectral entropy and frontal-temporal difference measures, which allowed the study to retain a compact interpretation layer alongside useful predictive signal.

We believe the article is well suited to \emph{Knowledge-Based Systems} because it addresses an artificial intelligence driven prediction and decision-support problem in healthcare, while combining methodological development with practical evaluation. The study contributes a sparse biomarker-discovery framework for low-channel EEG, examines how model capacity and sparse feature budgets affect discriminatory performance, and preserves interpretable biomarker structure rather than reporting predictive performance alone.

This submission is original, has not been published previously, is not under consideration elsewhere, and has been approved by all authors.

Thank you for your time and consideration.

Sincerely,

\vspace{1.2em}

Abin Daniel Zorto\\
Corresponding author\\
\texttt{u2091940@uel.ac.uk}

\end{document}
